% Generated draft: requires mathematical and source review.
\documentclass[11pt]{article}
\usepackage[margin=25mm]{geometry}
\usepackage{fontspec}
\setmainfont{Latin Modern Roman}
\usepackage{amsmath,amssymb,mathtools,unicode-math}
\setmathfont{Latin Modern Math}
\usepackage{enumitem,needspace,xurl,hyperref}
\setlength{\parindent}{0pt}
\setlength{\parskip}{5pt}
\setlength{\emergencystretch}{3em}
\newcommand{\N}{\mathbb{N}}
\newcommand{\Z}{\mathbb{Z}}
\newcommand{\Q}{\mathbb{Q}}
\newcommand{\R}{\mathbb{R}}
\newcommand{\C}{\mathbb{C}}
\newcommand{\F}{\mathbb{F}}
\newcommand{\A}{\mathbb{A}}
\newcommand{\PP}{\mathbb{P}}
\newcommand{\E}{\mathbb{E}}
\newcommand{\eps}{\varepsilon}
\newcommand{\dd}{\,\mathrm{d}}
\begin{document}
\section*{UnsolvedMath ID: 30004628}
\textbf{Source problem code:} catalogue entry
\subsection{Definitions and mathematical statement}
For a simple undirected graph \(\Gamma=(V,E)\), the associated right-angled Artin group is
\[
A_\Gamma=\left\langle v\ (v\in V)\ \middle|\ vw=wv\text{ whenever }\{v,w\}\in E\right\rangle.
\]
The supplied record does not specify whether the defining graph must be finite; that restriction remains unresolved. An automorphism \(\phi\in\operatorname{Aut}(A_\Gamma)\) is a bijective group homomorphism from \(A_\Gamma\) to itself. Its fixed subgroup is
\[
\operatorname{Fix}(\phi)=\{g\in A_\Gamma:\phi(g)=g\}\leq A_\Gamma.
\]
\textbf{Research question.} Given a right-angled Artin group \(A_\Gamma\) and an arbitrary automorphism \(\phi\in\operatorname{Aut}(A_\Gamma)\), describe the subgroup \(\operatorname{Fix}(\phi)\).

This preserves the supplied question, but it does not specify a particular property to establish or a proposed classification. Consequently, the record does not determine a unique conjectural assertion. No additional restriction on the automorphism is supplied.
\subsection{Short English statement}
Describe the subgroup of elements fixed by an arbitrary automorphism of a right-angled Artin group. The source leaves both the intended form of the description and the finiteness of the defining graph unspecified.
\subsection{Sources}
\begin{itemize}
\item Local UnsolvedMath registry, \texttt{data/unsolvedmath/problems.json}, permanent ID \(30004628\), source problem code \texttt{catalogue entry}, in the Oberwolfach Reports Open Problems collection. The supplied record lists proposed year \(2021\) and names no proposer.
\item Supplied report citation: \emph{Nonpositively Curved Complexes} (2022), Oberwolfach Reports, pp.\ 383--417, DOI \texttt{10.4171/owr/2021/8}. \url{https://doi.org/10.4171/owr/2021/8}. The background describes these reports as documenting mathematical workshops and their open-problem discussions.
\item Supplied literature reference: Jones--Vaskou, fixed subgroups in Artin groups. \url{https://arxiv.org/abs/2407.11839}.
\item Supplied literature reference: Lei--Zhang, fixed subgroups. \url{https://doi.org/10.1016/j.jalgebra.2023.07.006}.
\item Supplied literature reference: Oberwolfach report. \url{https://doi.org/10.4171/owr/2021/8}.
\item Literature assessment dated exactly \texttt{2026-08-22}: the supplied curation labels the problem open, states that fixed subgroups are understood for important subclasses while unrestricted automorphisms of general Artin or right-angled Artin groups remain unresolved, and says recent work explicitly poses the unrestricted large-type case as open. This is supplied dated curation derived from the OWR material, not an independently verified literature assessment. The record itself requests independent verification of time-sensitive status and cited claims.
\end{itemize}
\end{document}
